\section{Positive Laplace comparisons and fluctuations}\label{sec:fluctuations}

Throughout this section, $T=\beta^2>1$ is fixed.  We use the uniform
spin prior and write
\[
 H_n^{\mathrm{off}}(\sigma)
   =\frac{\beta}{\sqrt n}\sum_{i<j}g_{ij}\sigma_i\sigma_j,
 \qquad
 F_n=\log\left(2^{-n}\sum_{\sigma\in\{-1,1\}^n}
                         e^{H_n^{\mathrm{off}}(\sigma)}\right).
\]
Let $D$ be an independent centered Gaussian of variance $T/2$, and put
$F_n^c=F_n+D$.  Indeed,
\[
 \mathbb E H_n^{\mathrm{off}}(\sigma)H_n^{\mathrm{off}}(\tau)
       =\frac{nT}{2}R(\sigma,\tau)^2-\frac T2,
\]
so the completed Hamiltonian has covariance $nTR^2/2$, as required by
the definition of $f_n$.  Completion changes the variance by exactly
$T/2$, and independence gives the exact centered-transform correction
\[
 \log\mathbb E e^{t(F_n^c-\mathbb EF_n^c)}
 =\log\mathbb E e^{t(F_n-\mathbb EF_n)}+\frac{Tt^2}{4}.
\]
Changing the uniform prior to the counting measure adds the
deterministic constant $n\log2$ and has no effect on centered fluctuations.

Fix a constant $m$ with $0<m<1/300$, and set $\delta=m/10000$.
All assertions below hold for sufficiently large $n$, with constants
and thresholds allowed to depend on $T$ and $m$.  The order of choices
is important: $m$ is fixed before $n$ tends to infinity.  Define
\[
 s=n^{-1/6+5m},\qquad
 p(\varepsilon)=f_n(T-\varepsilon\mathbf1_{\{u<s\}},0),
 \qquad 0\le\varepsilon\le T.
\]
The endpoint with $\varepsilon=0$ integrates all matrix disorder at
mass zero, whereas the endpoint with $\varepsilon=T$ integrates it at
mass $s$.  The self correction is $T/4$ at both endpoints.  Consequently
\[
 p(0)=\frac1n\mathbb EF_n^c-\frac T4,\qquad
 p(T)=\frac1{ns}\log\mathbb E e^{sF_n^c}-\frac T4,
 \qquad
 K_n(s):=\log\mathbb E e^{s(F_n^c-\mathbb EF_n^c)}
               =ns\,[p(T)-p(0)].
\]
Our goal is the window in Proposition~\ref{fl:laplace-window}:
$c_{T,m}n^{20m}\le K_n(s)\le C_{T,m}n^{100m}$.
The lower and upper comparisons below establish its two bounds.

Moreover, \eqref{eq:1} gives
\[
 0\le p'(\varepsilon)=\frac14\int_0^s B_u\,du\le\frac s4
\]
at differentiability points, and the integrated inequality holds on
every interval.  In particular, $p$ is increasing.

\subsection{A chord inequality and the lower comparison}

We first record the form of the interpolation bound used throughout
this section.  Its hypotheses concern its two endpoint paths; a
pointwise difference of two admissible clocks need not itself be an
increasing clock.

\begin{lemma}[Chord comparison]\label{fl:chord}
Suppose $(b,h)$ and $(b+a,h-c)$ are bounded admissible paths,
$a\ge0$, and $c=0$ wherever $a=0$.  If $c^2/a$ is integrable, with
the quotient set to zero where $a=c=0$, then
\[
 f_n(b+a,h-c)\le f_n(b,h)+\frac14\int_0^1\frac{c(u)^2}{a(u)}\,du.
\]
\end{lemma}

\begin{proof}
The segment $(b+ta,h-tc)$ consists of admissible paths.  By
\eqref{eq:1}, its derivative is
\[
 -\frac14\int aB\,du+\frac12\int cS\,du
 \le\frac14\int(-aS^2+2cS)\,du
 \le\frac14\int\frac{c^2}{a}\,du,
\]
where $B\ge S^2$ and completion of the square give the two inequalities.
Integration proves the result.  The integrated variation identity
\eqref{eq:1}, established by step approximation, applies to the
bounded segments used here.
\end{proof}

Take $\varepsilon=n^{-1/3}$ and apply the short comparison
\eqref{eq:6} to the endpoint defining $p(\varepsilon)$, on all masses,
with
\[
 d=s,\qquad E=n^{-2/3+2m}d=n^{-5/6+7m}.
\]
Here $a=\sqrt n E/d=n^{-1/6+2m}$ and
$\varepsilon=o(a^2)$, so \eqref{eq:3} and \eqref{eq:4} hold.
All matrix paths stay uniformly positive.  The final-strength profile
has
\[
 e(r)=d^2f(r/d),\qquad
 n^\delta\frac{\sqrt n E}{d}=n^{-1/6+2m+\delta}=o(s).
\]
Thus \eqref{eq:8}, supplied by Proposition~\ref{reg:finite-propagation},
gives $r\asymp s$ on $s/2\le u\le s$ at a minimizing final path.
In particular $e(r)\asymp s^2\gg\varepsilon$ there.

Write $P$ for this changed path and $\mathcal P_e=\int er^2/4$ for
its penalty.  The comparison error gives
\[
 f_n(P)+\mathcal P_e\le p(\varepsilon)+n^{20\delta}E.
\]
Lemma~\ref{fl:chord}, with residual matrix covariance
$e+\varepsilon\mathbf1_{\{u<s\}}$ and residual field $er$, bounds
$p(0)$ above by the same changed pressure plus the corresponding
chord penalty.  Subtraction yields
\begin{equation}
\begin{split}
 p(\varepsilon)-p(0)
 &\ge\frac14\int_0^s
       \frac{\varepsilon e(r)}{e(r)+\varepsilon}\,r^2\,du
           -n^{20\delta}E\\
 &\ge c_{T,m}\varepsilon s^3.
\end{split}\tag{30}\label{eq:30}
\end{equation}
For the last step, the integral over $[s/2,s]$ is bounded below by a
constant times $\varepsilon s^3=n^{-5/6+15m}$, whereas the error is
$n^{-5/6+7m+20\delta}$.  Hence
\[
 K_n(s)\ge ns[p(\varepsilon)-p(0)]\ge c_{T,m}n^{20m}.
\]

\subsection{The initial upper comparison}

Set
\[
 d_H=n^{-1/6+15m},\qquad d_N=n^{-1/6+25m},\qquad
 \varepsilon_0=d_Hd_N=n^{-1/3+40m}.
\]
Apply the all-mass short comparison to $p(0)$ with $d=d_N$ and
$E=n^{-2/3+2m}d_N$.  At final strength, \eqref{eq:8} and monotonicity
give $r\le C d_N$ below mass $s$.  The fixed profile $f$ is positive
on bounded intervals, so $e(r)\ge c d_N^2$ there.  Since
$\varepsilon_0/d_N^2=n^{-10m}$, the residual matrix covariance
$e-\varepsilon_0\mathbf1_{\{u<s\}}$ is positive for large $n$.
Lemma~\ref{fl:chord} therefore gives
\[
\begin{split}
 p(\varepsilon_0)-p(0)
 &\le\frac14\int_0^s
       \frac{\varepsilon_0e(r)}{e(r)-\varepsilon_0}\,r^2\,du
       +n^{20\delta}E\\
 &\le C s\varepsilon_0d_N^2+n^{20\delta}E
 \le C_{T,m}\bigl(n^{-5/6+95m}
                    +n^{-5/6+27m+20\delta}\bigr).
\end{split}
\]
Only larger gaps in $\varepsilon$ remain.

\subsection{High comparisons with a low-field wall}

Fix $\varepsilon\in[\varepsilon_0,T)$ and put
\[
 q=\frac14\min(\varepsilon,T-\varepsilon),\qquad c_L=2q,
 \qquad d=d_H,\qquad E=n^{-2/3+10m}d.
\]
The high comparison acts on $A=[s,1]$.  Use the constraints
\eqref{eq:5}, with
\[
 u_*=(nE)^{-1},\qquad r_{\min}=\frac{u_*}{2M(s)},\qquad
 r_0=\frac{d^2r_{\min}}{2\varepsilon}.
\]
Its final strength is $e(r)=d^2f(r/d)$.  To specify the low field,
let $\tau=n^{-10}$ and choose
\[
 \psi(u)=\frac{\log((s+\tau)/(s-u+\tau))}
                    {\log((s+\tau)/\tau)},\qquad 0\le u\le s.
\]
Thus $\psi(0)=0$, $\psi(s)=1$, $\psi$ is increasing, and
\[
 \psi'(u)\ge\frac{c}{(s-u)\log n}
               \quad\text{when }s-u\ge\tau.
\]

For $0\le\theta\le1$, minimize over the high quantile $r$ the
functional
\[
\begin{split}
 J_\theta(r)={}&
 f_n\bigl(T-\varepsilon\mathbf1_{\{u<s\}}
            -\theta c_L\mathbf1_{\{u<s\}}-e(r)\mathbf1_A,\,
          \theta c_Lr_0\psi\mathbf1_{\{u<s\}}
                              +e(r)r\mathbf1_A\bigr)\\
 &+\frac14\int_Ae(r)r^2\,du
       +\frac{\theta c_L}{4}\int_0^s r_0^2\psi(u)^2\,du.
\end{split}
\]
Write $J_\theta^*$ for its minimum.  All references below to
$S,B,D,K$ for this functional refer to an actual minimizing changed
path, not to the unmodified low-field coefficients.

\begin{lemma}[Admissibility and the unswept wall profile]
\label{fl:wall-profile}
For all sufficiently large $n$, uniformly in
$\varepsilon\in[\varepsilon_0,T)$ and $\theta\in[0,1]$, every path
in the definition of $J_\theta$ is admissible.  Its minimizing high
quantiles have the regular outer profile \eqref{eq:8}; in particular
they are regular above scale $d$.  In using this profile assertion,
no pressure is evaluated at the bare low-field path with its high
comparison removed.
\end{lemma}

\begin{proof}
Direct substitution in \eqref{eq:5} gives
\[
 a=n^{-1/6+10m},\quad u_*=L=n^{-1/6-25m},\quad
 M(s)\asymp n^{30m+\delta},\quad
 r_{\min}\asymp n^{-1/6-55m-\delta}.
\]
In particular, $s=o(a)$, $d^2=o(\varepsilon_0)$, and
$r_{\min}/d\to0$.  Since $f=1$ near zero, the initial high field is
exactly $d^2r_{\min}$.  On the low group the field is at most
\[
 c_Lr_0\le\frac\varepsilon2
                   \frac{d^2r_{\min}}{2\varepsilon}
             =\frac{d^2r_{\min}}4.
\]
There is therefore no downward field jump at $s$.  Both field paths
are increasing on their own intervals.  The low matrix coefficient
is at least $(T-\varepsilon)/2\ge0$, and its gap to the high coefficient
at the wall is $\varepsilon+\theta c_L-d^2>0$.  The high matrix path
is increasing because $e$ decreases.

These are the wall data of Corollary~\ref{reg:wall-uniformity}, with
$\varphi=\psi$. Total clocks are uniformly bounded, and $d^2<T/2$
for large $n$. The fixed covariance margin required there is also
uniform: when $\varepsilon<T/4$, the low covariance is at least
$5T/8$ and the high covariance is at least $T-d^2$; when
$\varepsilon\ge T/4$, the high-to-low gap is at least $T/4-d^2$.
These are respectively the all-mass and high-only alternatives in
the corollary. Above $s$, the bare coefficients are exactly $T,0$,
so \eqref{eq:4} holds uniformly.

The low penalty in $J_\theta$ is independent of the high quantile.
It therefore leaves its minimizers unchanged and cancels in the
movement inequality used by the corollary. Applying the corollary
at this fixed positive strength gives the profile uniformly along
all wall-parameter sequences, including $\theta\to0$ and
$T-\varepsilon\to0$, without evaluating the bare low-field pressure
or sweeping the root from zero strength. Here its final-strength
cutoff is
\[
 A_{d^2}=n^\delta a=n^{-1/6+10m+\delta}=o(d).
\]
Thus the resulting outer regularity applies at scale $d$, as claimed.
\end{proof}

\begin{lemma}[Suppression of the low-group mean]\label{fl:low-mean}
For every minimizing changed path in Lemma~\ref{fl:wall-profile},
uniformly over its low-group splits,
\begin{equation}
 S_u=\frac{K_u}{T}+O(d^3),\qquad
 S_u\le\frac{C d^3}{\varepsilon}=:a_L,
 \qquad K_u=h_u+b_uS_u,\quad 0\le u<s.
 \tag{31}\label{eq:31}
\end{equation}
The constants are uniform in the wall parameters.
\end{lemma}

\begin{proof}
At scale $x=d$, the final strength satisfies $e(d)\asymp d^2$, so
$W\asymp d$. The cutoff just proved supplies outer regularity there,
and the other direct-test inequalities follow from
\[
 d/a=n^{5m},\qquad
 \frac{nE^2}{e(d)d^2}\asymp n^{-10m}.
\]
Also $ne(d)d^4\asymp n^{90m}\to\infty$.
Projection from a regular high
band and \eqref{eq:12} give $K_u=O(d)$ and $D_u=O(d)$ on the omitted
low group.  The omitted-group big-O version of \eqref{eq:22} therefore
gives $S_u=\Gamma(K_u)+O(d^3)$.  Its hypotheses require an admissible
changed path and an omitted mass $s=o(d)$; they do not require
\eqref{eq:4} at the omitted splits.

Choose a fixed $C$ containing all low-group clocks in $[0,Cd]$.
Lemma~\ref{prof:low-clock-linearization}, applied at $x=d$ with
omitted mass $u_0=s=o(d)$, gives
\[
 \sup_{0\le K\le Cd}|\Gamma(K)-K/T|=O(d^3)
\]
uniformly in the wall parameters. Its high-clock calibration
covers this interval down to zero, including clock gaps, and
requires no version of \eqref{eq:4} on the omitted group.
Combining this bound with the preceding omitted-group cavity
estimate proves the first assertion of \eqref{eq:31}.

On the low group,
$b_u=T-\varepsilon-\theta c_L$ and
$0\le h_u\le d^2r_{\min}/4=O(d^3)$.  Therefore
\[
 (\varepsilon+\theta c_L)S_u=h_u+O(d^3)\le C d^3,
\]
which proves the second assertion.  Uniformity follows either from
the uniform constants just used or by applying the same argument to
an arbitrary sequence of wall parameters and low splits.
\end{proof}

\subsection{The low-field sweep and the upper Laplace bound}

\begin{lemma}[Cost of a wall sweep]\label{fl:wall-cost}
With the preceding choices, uniformly in
$\varepsilon_0\le\varepsilon<T$,
\[
\begin{split}
 J_1^*-J_0^*
 \le C\left[
    c_Ls(a_L^2+r_0^2)
       +\frac{a_L(\log n)^3}{nr_0s}+n^{-10}\right].
\end{split}
\]
Consequently,
\[
 p(\varepsilon+q)-p(\varepsilon)
 \le n^{20\delta}E
       +C\left[s\frac{d^6}{\varepsilon_0}
              +\frac{d}{nsr_{\min}}(\log n)^3+n^{-10}\right].
\]
\end{lemma}

\begin{proof}
We first take $\theta\ge\tau=n^{-10}$. For the dyadic values
$v=s/2,s/4,\ldots$ with $v\ge4\tau$, put
\[
 I_v^-=[s-v,s-v/2],\qquad I_v^+=[s-v/2,s-v/4].
\]
The intervals $I_v^-$ cover $[s/2,s)$ except for a terminal strip of
length at most $4\tau$. Each $I_v^+$ lies after $I_v^-$, has length
$v/4$, and stays at distance at least $\tau$ from $s$. On either
interval, the field derivative is at least
$c\theta c_Lr_0/(v\log n)$. Integrating the susceptibility
inequality in \eqref{eq:1} and using \eqref{eq:31} gives
\[
 \int_J\mathbb E\operatorname{tr}H_u^2\,du
          \le\frac{C a_Lv\log n}{n\theta c_Lr_0}.
\]
Here $J$ is either $I_v^-$ or $I_v^+$. Since $H_u\ge u C_u$ and
$u\ge s/2$ on both, it also gives
\[
 \int_J\mathbb E\operatorname{tr}C_u^2\,du
          \le\frac{C a_Lv\log n}{n\theta c_Lr_0s^2}.
\]

For $u\in I_v^-$, average the single-path observable over $I_v^+$:
\[
 \overline L_v=\frac1{|I_v^+|}\int_{I_v^+}L_j\,dj,\qquad
 |X_u|^2\le2\mathbb E_u\overline L_v.
\]
Indeed $2\mathbb E_u L_j=\mathbb E_u|X_j|^2\ge|X_u|^2$ for $j\ge u$.
Moreover $\mathbb E\overline L_v\le a_L/2$. Apply the pinned moment inequality
in \eqref{eq:1} with $p=2$ and the deterministic weight
$k_j=|I_v^+|^{-1}\mathbf1_{I_v^+}(j)$. Since
\[
 \mathbb E W_j=\mathbb E\operatorname{tr}(H_jC_j)
                    \le\frac2s\mathbb E\operatorname{tr}H_j^2,
\]
the variance under the full planted law satisfies
\[
 \operatorname{Var}(\overline L_v)
 \le\frac1{|I_v^+|^2}\int_{I_v^+}\mathbb E W_j\,dj
 \le\frac{C a_L\log n}{n\theta c_Lr_0sv}.
\]
Conditional Jensen is used only after this unconditional estimate:
\[
 \mathbb E|X_u|^4
 \le4\mathbb E\overline L_v^2
 \le C a_L^2+
             \frac{C a_L\log n}{n\theta c_Lr_0sv}.
\]
Also
\[
 B_u=\mathbb E\|X_uX_u^{\mathsf T}+C_u\|_{\mathrm{HS}}^2
 \le2\mathbb E|X_u|^4+2\mathbb E\operatorname{tr}C_u^2.
\]
On integrating over $I_v^-$, the fourth-moment error loses its factor
$v^{-1}$. The covariance term is no larger, because $v\le s$:
\[
 \int_{I_v^-}B_u\,du
 \le Cv a_L^2+
      \frac{C a_L\log n}{n\theta c_Lr_0s}.
\]
Summing the $O(\log n)$ dyads proves
\[
 \int_{s/2}^{s-C\tau}B_u\,du
 \le Csa_L^2+
                \frac{C a_L(\log n)^2}{n\theta c_Lr_0s}.
\]
The omitted terminal strip has cost at most $C\tau$, since $B_u\le1$.
For $u<s/2$, monotonicity gives
$\int_0^{s/2}B_u\,du\le2\int_{s/2}^{3s/4}B_u\,du$.
Hence the same bound, with a $C\tau$ remainder, holds over
the whole low group.

The feasible high-quantile class is fixed, so the pressure Lipschitz
bound makes $J_\theta^*$ Lipschitz in $\theta$.
At an optimizing high quantile, the upper envelope derivative per
unit $\log\theta$ is at most
\[
 \frac{\theta c_L}{4}\int_0^s
       [D_u^2+(S_u-r_0\psi(u))^2]\,du.
\]
This is \eqref{eq:7} for the low variation, with the high optimizer
held fixed when taking an upper difference quotient.  The estimate
is uniform over all minimizers, so it also bounds the derivative of
the minimum at almost every parameter.  Dropping the nonpositive
cross term and using $D_u^2+S_u^2=B_u$ bounds it by
\[
 C\theta c_Ls(a_L^2+r_0^2)
       +\frac{C a_L(\log n)^2}{nr_0s}+C\theta c_L\tau.
\]
Integration over $\tau\le\theta\le1$ gives the claimed three-logarithm
term.  On $0\le\theta\le\tau$ the trivial bounds on the overlap and
the fixed low coefficient give total cost $C\tau c_Ls\le C\tau$.
This proves the first assertion.

At $\theta=0$, the high comparison begins from the admissible endpoint
defining $p(\varepsilon)$ and can be swept from strength zero.  Hence
$J_0^*\le p(\varepsilon)+n^{20\delta}E$.
At $\theta=1$, the chord to $p(\varepsilon+q)$ has residual matrix
covariance $e$ on the high group and $c_L-q=c_L/2$ on the low group.
Its residual field is $er$ high and $c_Lr_0\psi$ low.  Relative to
the penalties included in $J_1^*$, the excess chord cost is
\[
 \frac14\int_0^s
  \left(\frac{c_L^2}{c_L-q}-c_L\right)r_0^2\psi^2\,du
       =\frac{c_L}{4}\int_0^s r_0^2\psi^2\,du
       \le Csc_Lr_0^2.
\]
Both endpoints and their connecting segment are admissible.  Finally,
$a_L=Cd^3/\varepsilon$, $c_L\le\varepsilon/2$, and
$r_0/a_L\le C r_{\min}/d\to0$.  These imply
\[
 c_Ls(a_L^2+r_0^2)\le Csd^6/\varepsilon_0,
 \qquad \frac{a_L}{r_0}\le C d/r_{\min},
\]
and prove the second assertion.
\end{proof}

Here is the complete power ledger for these choices; constants and
the indicated logarithms are suppressed in the table.
\begin{center}
\begin{tabular}{ll}
\toprule
Quantity & Power of $n$\\
\midrule
$\varepsilon s^3$ in the lower comparison & $-5/6+15m$\\
Lower comparison error $n^{20\delta}E$ & $-5/6+7m+20\delta$\\
Initial upper chord $s\varepsilon_0d_N^2$ & $-5/6+95m$\\
Initial upper comparison error & $-5/6+27m+20\delta$\\
High comparison error $n^{20\delta}E$ & $-5/6+25m+20\delta$\\
$sd^6/\varepsilon_0$ & $-5/6+55m$\\
$d/(nsr_{\min})$ & $-5/6+65m+\delta$\\
\bottomrule
\end{tabular}
\end{center}
For reference, the ceiling condition in \eqref{eq:3} has respective
ratios $E/(n^{-1/2}d^2)=n^{-3m},n^{-23m},n^{-5m}$ in the lower,
initial upper, and high comparisons.  Its lower condition follows
from the respective factors $n^{2m},n^{2m},n^{10m}$ in the choices
of $E/d$.  Thus none of these applications sits at an unverified
endpoint of the allowable error range.

Starting at $\varepsilon_0$, iterate
$\varepsilon\mapsto\varepsilon+\frac14\min(\varepsilon,T-\varepsilon)$
until $T-\varepsilon\le n^{-5}$.  Before crossing $T/2$, the current
$\varepsilon$ grows by a factor $5/4$; subsequently its distance to
$T$ decreases by a factor $3/4$.  The number of steps is therefore
$O_{T,m}(\log n)$.  The last unswept interval costs at most $sn^{-5}/4$.
The ledger and Lemma~\ref{fl:wall-cost}, uniformly at every step, imply
\[
 p(T)-p(0)\le C_{T,m}\left[
 n^{-5/6+95m}
 +n^{-5/6+65m+\delta}(\log n)^4
 +n^{-5/6+27m+20\delta}\log n+n^{-5}s\right].
\]
Every term after the first is smaller for each fixed $m>0$ and
sufficiently large $n$.  We have proved the following moment window.

\Needspace{3\baselineskip}
\begin{proposition}[A positive Laplace window]\label{fl:laplace-window}
For every fixed $T>1$ and $0<m<1/300$, there exist $c_{T,m}>0$,
$C_{T,m}<\infty$, and $n_0(T,m)$ such that
\[
 c_{T,m}n^{20m}\le
 \log\mathbb E\exp\{s(F_n^c-\mathbb EF_n^c)\}
 \le C_{T,m}n^{100m},\qquad s=n^{-1/6+5m},
\]
for all $n\ge n_0(T,m)$.
\end{proposition}

\subsection{Exponential augmentation and one-dimensional estimates}

A positive moment by itself need not determine a variance: an
exponentially rare positive deviation can dominate it, and a rare
negative deviation can dominate the variance without enlarging that
moment.  The next structural property supplies the needed distributional
control.  It is the exponential augmentation underlying the positive
moment method of \cite[Proposition~5.1]{AL2026}; we give its proof in
the present normalization.

\begin{lemma}[Log-concave augmentation]\label{fl:augmentation}
Let $\mathcal E$ be an independent exponential variable of rate one.
Then $G_n=F_n^c+\mathcal E$ has a log-concave density.  For $0\le t<1$,
\[
 \log\mathbb E e^{t(G_n-\mathbb EG_n)}
       =K_n(t)-t-\log(1-t),\qquad
 \operatorname{Var}G_n=\operatorname{Var}F_n^c+1.
\]
\end{lemma}

\begin{proof}
Represent all disorder, including the common Gaussian, by a standard
Gaussian vector $g\in\mathbb R^N$.  There are linear forms
$a_\sigma\cdot g$ such that
\[
 F_n^c(g)=\log\left(2^{-n}\sum_\sigma e^{a_\sigma\cdot g}\right).
\]
For each $\tau\in\{-1,1\}^n$, the orthogonal gauge transformation
$g_{ij}\mapsto\tau_i\tau_jg_{ij}$, leaving the common coordinate fixed,
preserves both the Gaussian law and $F_n^c$.  These transformations
act transitively on the spin labels.  The density of $G_n$ is
\[
 q(z)=e^{-z}\mathbb E\left[e^{F_n^c(g)}
                           \mathbf1_{\{F_n^c(g)\le z\}}\right].
\]
Upon expanding $e^{F_n^c}$, gauge invariance shows that all terms
have the same integral.  Thus, for one fixed spin label $\sigma_*$,
$q$ is the marginal, up to a constant, of
\[
 \exp\{-|g|^2/2+a_{\sigma_*}\cdot g-z\}
                      \mathbf1_{\{F_n^c(g)\le z\}}.
\]
The function $F_n^c$ is convex, being a logarithm of a sum of
exponentials of linear functions.  Its epigraph is convex, and the
displayed integrand is therefore log-concave jointly in $(g,z)$.
It is integrable: integrating first in $z$ leaves a Gaussian density
times $e^{a_{\sigma_*}\cdot g-F_n^c(g)}\le2^n$.
Pr\'ekopa's marginal theorem \cite[Theorems~6 and~8]{Prekopa1973}
implies that $q$ is log-concave. Independence and
$\mathbb E e^{t(\mathcal E-1)}=e^{-t}/(1-t)$ give the two identities.
\end{proof}

\begin{lemma}[Variance and concentration for a log-concave law]
\label{fl:one-dimensional}
Let $X$ be centered with a log-concave density and
$\sigma^2=\mathbb E X^2>0$.  There are universal positive constants
$c,C$ such that
\[
\begin{split}
 &\mathbb P\{|X|\ge t\sigma\}\le2e^{-ct}\quad(t\ge0),\\
 &\log\mathbb E e^{tX}\le C t^2\sigma^2
                       \quad(0\le t\sigma\le c),\\
 &\mathbb P\{X\ge c\sigma\}\ge c,\qquad
   \log\mathbb E e^{tX}\ge ct\sigma+\log c\quad(t\ge0),\\
 &\sup_x f_X(x)\le C/\sigma.
\end{split}
\]
The lower log-moment inequality is understood also when the moment is
infinite.
\end{lemma}

\begin{proof}
The survival functions of $X$ and $-X$ are log-concave by
Pr\'ekopa's theorem.  Chebyshev's inequality bounds the survival
probability at $-2\sigma$ below by $3/4$, and that at $2\sigma$ above
by $1/4$.  Concavity of its logarithm and the secant slope between
these two points give, for $t\ge2$,
\[
 \mathbb P(X\ge t\sigma)\le\frac{\sqrt3}{4}
                    \exp\left(-\frac{\log3}{4}t\right).
\]
Apply the same argument to $-X$; the trivial bound for $t\le2$
then gives the first assertion with $h=(\log3)/4$ in place of $c$.
Integration of the tails yields
\[
 \mathbb E|X|^k\le2k!(\sigma/h)^k,\qquad k\ge1.
\]
For $r=t\sigma/h\le1/2$, the exponential series converges absolutely
and, since $\mathbb EX=0$,
\[
 \log\mathbb E e^{tX}
 \le\log\left(1+2\sum_{k\ge2}r^k\right)
 \le\frac{2r^2}{1-r}\le4r^2.
\]

The fourth-moment bound and H\"older's inequality give
\[
 \mathbb E|X|\ge
       \frac{(\mathbb EX^2)^{3/2}}{(\mathbb EX^4)^{1/2}}
       \ge c_1\sigma.
\]
Centering implies $\mathbb EX_+=\mathbb E|X|/2$, whereas
$\mathbb EX_+^2\le\sigma^2$.  The Paley--Zygmund inequality, or its
elementary proof by Cauchy--Schwarz on the event
$\{X_+\ge\mathbb EX_+/2\}$, now gives
$\mathbb P(X\ge c_2\sigma)\ge c_2$ for a universal $c_2>0$.
Restricting the exponential moment to that event proves its lower
bound.

For the density estimate, write $M=\sup f_X$ and choose a mode $x_0$.
An integrable one-dimensional log-concave density has a finite positive
maximum, with one-sided limits used at a support endpoint.  On either
side of the mode the density must fall to $M/e$ within distance at
most $e/M$, or reach the support endpoint: otherwise its integral on
that side would exceed one.  Concavity of the log density then bounds
the tail beyond that point by its exponentially decreasing secant.
Together with $f_X\le M$ before that point, this gives
\[
 f_X(x_0+y)\le eM\exp(-M|y|/e),
\]
with zero outside the support.  Consequently
\[
 \sigma^2\le\mathbb E(X-x_0)^2
       \le2eM\int_0^\infty y^2e^{-My/e}\,dy
       =\frac{4e^4}{M^2}.
\]
Rearranging proves $M\le2e^2/\sigma$.
\end{proof}

\subsection{Conclusion with fixed power slack}

\begin{proof}[Proof of Theorem~\ref{thm:main}]
Let $\sigma_n^2=\operatorname{Var}G_n$.  For the fixed $m$ used in
Proposition~\ref{fl:laplace-window}, $s\to0$.  The exponential
correction in Lemma~\ref{fl:augmentation} is $O(s^2)$, and hence
\[
 c_{T,m}n^{20m}\le
 \log\mathbb E e^{s(G_n-\mathbb EG_n)}
                         \le C_{T,m}n^{100m}
\]
for sufficiently large $n$, after changing constants.  If
$s\sigma_n$ were smaller than the universal small-tilt constant
in Lemma~\ref{fl:one-dimensional}, its upper quadratic bound would
make this log moment uniformly bounded.  The divergent lower bound
therefore gives $s\sigma_n\ge c$ eventually.  Its positive-side mass
bound, on the other hand, gives
\[
 cs\sigma_n+\log c
       \le\log\mathbb E e^{s(G_n-\mathbb EG_n)}
       \le C_{T,m}n^{100m}.
\]
It follows that
\[
 c_{T,m}n^{1/6-5m}\le\sigma_n
                  \le C_{T,m}n^{1/6+95m}.
\]
Given any desired power slack $\eta>0$, first choose a fixed
$m<\min(1/300,\eta/200)$.  The constants are absorbed by the remaining
strict power slack as $n\to\infty$.  Thus
$\sigma_n=n^{1/6+o(1)}$.  Since
\[
 \operatorname{Var}G_n=\operatorname{Var}F_n+T/2+1,
\]
the same standard-deviation exponent holds for $F_n$, and its variance
is $n^{1/3+o(1)}$.

We also prove the asserted typical magnitude.  Let
$X_n=G_n-\mathbb EG_n$.  The density estimate in
Lemma~\ref{fl:one-dimensional} implies, for every $a>0$,
\[
 \mathbb P(|X_n|\le a)\le 2Ca/\sigma_n.
\]
Using the standard-deviation bounds with any fixed slack smaller
than $\eta$, this tends to zero at $a=n^{1/6-\eta}$.  Chebyshev's
inequality, or the exponential-tail bound of that lemma, shows that
$\mathbb P(|X_n|\ge n^{1/6+\eta})\to0$.  Hence the smoothed centered
magnitude lies between these two powers with probability tending
to one.

To transfer this conclusion, write
\[
 X_n=(F_n-\mathbb EF_n)+N,\qquad
 N=D+(\mathcal E-1),\qquad \mathbb EN^2=T/2+1.
\]
First suppose $0<\eta<1/6$.  At the lower threshold
$a_n=n^{1/6-\eta}$, the event $|F_n-\mathbb EF_n|\le a_n$ is
contained in $\{|X_n|\le2a_n\}\cup\{|N|>a_n\}$; both probabilities
tend to zero.  At the upper threshold $b_n=n^{1/6+\eta}$, the event
$|F_n-\mathbb EF_n|>b_n$ is contained in
$\{|X_n|>b_n/2\}\cup\{|N|>b_n/2\}$, again with vanishing probability.
Fixed numerical factors are harmless because the standard-deviation
estimates are available with arbitrarily smaller fixed slack.
For larger $\eta$, use inclusion from any smaller positive slack
less than $1/6$.  This proves the typical-magnitude assertion for
every fixed $\eta>0$.

Finally, division by $n$ changes the standard-deviation and typical
magnitude exponents to $-5/6$.  The physical free-energy convention
introduces the fixed nonzero factor $-1/\beta$, and the counting spin
prior introduces only a deterministic constant.  These changes do
not alter either conclusion.  All constants and thresholds may depend
on the fixed $T>1$; no assertion uniform at either temperature endpoint
is used or obtained.
\end{proof}
